% Generated from reports/FIGURE_LEGENDS.md by scripts/63_legends_to_tex.py
% Do not edit here; edit the Markdown and regenerate.

\newcommand{\legendfigone}{%
\textbf{Acoustic arousal models transfer between corpora of the same sound domain but not across it, under every representation tested.} \textbf{a} Generalisation stages for six regression algorithms (RF, XGBoost, GBDT, SVR, KNN, ElasticNet) over 122 spectro-temporal descriptors. Each line follows one algorithm from within-corpus cross-validation, through same-domain transfer (music $\leftrightarrow$ music), to cross-domain transfer (music $\leftrightarrow$ environmental); grey points are the individual corpus pairs behind each mean. Within-corpus values are GroupKFold-5 estimates grouped by source recording; transfer values fit on the whole of one corpus and are evaluated on the whole of another, with nothing refitted. \textbf{b} Cross-corpus transfer matrix, averaged over the six algorithms; the diagonal shows within-corpus cross-validation. Green indicates positive rank correlation, red negative. \textbf{c} Directional asymmetry. For each music corpus, transfer outward to the environmental corpus and inward from it, averaged over algorithms; the dashed line marks the same-domain mean (0.59). The inward direction is uniformly sealed ($\rho$ $\le$ 0.09 for all three corpora), whereas the outward direction is not and increases from PMEmo ($-$0.04) through DEAM (+0.13) to Soundtracks (+0.24). Film soundtracks contain substantial atmospheric material --- sustained tones, texture, sound design --- and are correspondingly closer to soundscape recordings. The barrier is therefore \textbf{not symmetric}: its density in the outward direction tracks the composition of the source corpus, while the inward direction is impermeable regardless of target. \textbf{d} Duration control. Music clips truncated from 30 s to a 6 s centred window, matching the environmental corpus; filled circles are 30 s, open circles 6 s. Same-domain transfer survives truncation (0.67 $\rightarrow$ 0.63) while cross-domain transfer stays at zero (0.03 $\rightarrow$ 0.05), excluding clip length as the cause of the collapse. \textbf{e} Leave-one-category-out generalisation inside the environmental corpus (n = 1,213), holding out each of the six Schafer soundscape categories in turn. The shaded band marks the mean cross-domain level. Generalisation to unseen content categories holds (mean $\rho$ = 0.52; range 0.20 for \texttt{quiet} to 0.70 for \texttt{human}, five of six categories above 0.33), so the collapse in \textbf{a}--\textbf{b} is specific to the domain border rather than to any content shift. \textbf{f} The same operation costed twice, under five representations. For each representation the corpus is swapped and nothing is refitted; the grey mark gives the performance lost when the new corpus is in the same domain, the red mark the loss when it is not, both as a fraction of that representation's own within-corpus performance, and the printed value is the gap between them in percentage points. Points are medians over the corpus pairs of \textbf{a}. For the three uncontaminated families the loss is 20--31\% within the domain and 84--95\% across it --- hand-specified 20.2/93.9\%, AST 22.0/84.2\%, MERT 29.1/94.9\%, Wav2Vec2 30.7/90.0\%, a gap of 59 to 74 percentage points (Table 1). The result is where the collapse stands, not how high it is: every representation transfers between corpora and none crosses between domains, which disposes of the reading in which the embeddings simply fail to move between datasets. AST is the family that most ought to cross, since AudioSet contains large numbers of environmental categories; it raises cross-domain $\rho$ from 0.047 to 0.125, a factor of 2.7 that is still a sixth of what the same representation reaches within a corpus. CLAP is drawn apart and greyed. Its pretraining set contains Freesound, the source of the environmental corpus, an overlap declared before any of these numbers existed. It is not a generally stronger representation --- 0.620 on music-to-music against AST's 0.619 and 0.614 for the hand-specified descriptors --- and its lower barrier (57.3\%) appears only in comparisons that necessarily involve the corpus its pretraining data overlap. With a single environmental corpus available, genuine crossing cannot be separated from familiarity with the target; the barrier claim rests on the three uncontaminated families, and CLAP is used neither as support nor as refutation. Corpora: DEAM (music, n = 1,233), PMEmo (music, n = 717), Soundtracks (music, n = 360) and Emo-Soundscapes (environmental, n = 1,213). Soundtracks (Eerola \& Vuoskoski, 2011) differs from the other music corpora in collection site, decade, rating protocol (nine-point Likert scales over eight emotion terms rather than continuous valence--arousal sliders) and material selection (a balanced design of twelve target emotions $\times$ 30 excerpts). Its inclusion therefore tests whether within-music transfer reflects shared acoustic structure or merely a shared annotation tradition; transfer into it from the other music corpora is 0.56, and from the environmental corpus 0.09. Panels \textbf{a}--\textbf{e} rest on 12 ordered corpus pairs, each evaluated with the six algorithms: over the 36 same-domain values transfer averages 0.592 (s.d. 0.130) and over the 36 cross-domain values 0.054 (s.d. 0.161). Target is continuous perceived arousal rescaled to $-$1\ldots{}1. All audio is loudness-normalised to $-$23 LUFS before feature extraction, which is an admission criterion for this figure rather than a cosmetic step (Extended Data Fig. 1). Metric is Spearman $\rho$ throughout. Source data are provided as a Source Data file.}

\newcommand{\legendfigtwo}{%
\textbf{Target-side labels buy back most of the corpus-swap gap, and the purchase stops well short of the ceiling.} \textbf{a} Fraction of the gap recovered by k labelled target-corpus examples. The gap is measured between zero-shot transfer (k = 0, defined as 0\%) and a model trained on the target corpus itself (100\%), so the denominator is the target's own within-corpus ceiling and not same-domain transfer --- the question an engineer faces is how close labelling k examples gets to labelling all of them. Line and points, median over the six ordered corpus pairs; band, interquartile range across those pairs. Ten labels recover 20\% of the gap, 25 recover 34\%, 50 recover 50\%, 100 recover 64\% and 200 recover 69\%. \textbf{b} The same result in marginal terms: additional gap recovered per doubling of the budget. Each doubling up to 100 labels returns a roughly constant 14 to 16 percentage points --- the signature of a gap that is genuinely for sale --- while the doubling from 100 to 200 returns 5. The last bar is the point of the panel: the curve flattens at 69\% of the gap, so about a third of it is unbought at the largest budget tested, and this experiment does not say whether the remainder is purchasable at a budget we did not test or is a component of a different kind. \textbf{c} Alignment without labels, with its sham control. Correlation alignment (CORAL) recovers 22\% of the gap. Aligning the source not to the target but to an unrelated third corpus --- a manipulation that cannot carry target-specific information by construction --- recovers 14\%, so roughly two-thirds of the apparent gain is not target-specific and is whatever a second-moment correction does for any distribution. Diagonal CORAL recovers 5\% and per-corpus z-scoring 6\%. Subspace alignment is worse than doing nothing ($-$5\%); the axis is not clipped at zero, because clipping would turn ``actively harmful'' into ``small''. Every fraction is computed within an ordered corpus pair and only then pooled. Pooling first --- averaging the correlations and taking one ratio against one median ceiling --- returns 36\% and 27\% in place of the 22\% and 14\% in \textbf{c}, because it mixes pairs whose baselines differ by more than the effect being measured. The six pairs are the three music corpora against the environmental corpus in both directions, so the quantity priced here is the crossing of the music/environmental boundary. The estimator is ridge regression on the hand-specified descriptors; tree ensembles are reported in the Supplementary Information, and are kept separate because they differ in sensitivity to the alignment step. The 200-label point rests on five of the six pairs: Soundtracks (n = 360) cannot give up 200 labelled examples and still be evaluated. Metric is Spearman $\rho$; all audio is loudness-normalised to $-$23 LUFS. Source data are provided as a Source Data file.}

\newcommand{\legendfigthree}{%
\textbf{Feature-selection routes recover different descriptor classes depending on what they test.} \textbf{a} UpSet-style recovery pattern. The matrix marks which of four selection routes recovered each descriptor; connected points indicate a descriptor recovered by more than one route. The top bars count routes per descriptor and the left bars give each route's set size, split by feature class. The four routes are: SHAP importance with per-feature permutation FDR (q = 0.10, 200 permutations); preservation of the sign of the cross-domain Cohen's d; native shape-function importance from an Explainable Boosting Machine; and stability across a Rashomon set of 10 near-optimal models sampled from 60 ($\rho$ $\ge$ $\rho$\_max $-$ 0.02). \textbf{b} Share of each route's recovered set that consists of temporal-variability descriptors, against the 50\% chance line. The route that tests within-domain statistical significance returns 2 of 12 (17\%), below chance; the three routes that probe generalisation or model multiplicity return 5 of 8 (63\%), 6 of 10 (60\%) and 4 of 6 (67\%). The four routes differ in statistical principle --- post-hoc attribution, effect size, additive shape functions, and model multiplicity --- not in hyper-parameters, so \textbf{b} is a dissociation and not a shortfall of tuning. The route that is standard in the field returns, at a rate below chance, precisely the spectral-level descriptors that Fig. 1 and Extended Data Fig. 8 show do not survive the domain boundary. What the panel licenses is therefore a negative: no attribution method here supports a statement about what will transfer, which is why every boundary in this work is crossed and measured rather than predicted from the descriptors a model relies on inside its own corpus. The companion requirement --- that an apparent overlap between attribution sets be read against the overlap that shuffled labels produce --- is in Fig. 6a,b. Source data are provided as a Source Data file.}

\newcommand{\legendfigfour}{%
\textbf{A pre-declared falsification test: an information-limited gap that a change of representation does close.} \textbf{a} Cross-validated Spearman $\rho$ for eight affect axes of the Soundtracks corpus (n = 360), using 122 spectro-temporal descriptors alone (open circles) and with 39 tonal descriptors added (filled). Axes are ordered by gain; numbers give the gain. Blue marks the three axes for which musical mode is theoretically decisive (happiness, valence, sadness) --- a grouping fixed a priori, not read off the result. The criterion was fixed before the test, because adding 39 descriptors will improve something somewhere: the gain had to concentrate on happiness, at more than twice the mean gain of the other seven axes and above 0.05 in absolute terms. It does, at +0.132 against a mean of +0.045, a ratio of 2.9. The ordering of gains follows mode-dependence: happiness, valence (+0.083) and sadness (+0.077) lead, while energy (+0.012) and anger (+0.025), which depend on level, roughness and rhythm, gain least. \textbf{b} Tonal descriptors alone against spectro-temporal descriptors alone. For happiness, 39 tonal descriptors ($\rho$ = 0.580) outperform 122 spectro-temporal ones (0.474); the deficit was one of representation rather than of dimensionality. \textbf{c} Mechanism check. Rank of five tonal descriptors among all 161, from random forest importance, for the happiness and tension axes (log scale). The two families dissociate: major-key strength and mode score rank 1st and 2nd for happiness but 3rd and 20th for tension, whereas interval consonance and sensory roughness rank 4th and 7th for tension but 117th and 76th for happiness. Mode carries perceived happiness; consonance and roughness carry perceived tension. The ordering was selected by the model from 161 descriptors, not imposed. This figure is reported at the length it is because it is a pre-declared test of the distinction that organises this work, not a minor improvement. The condition was stated in advance: the distinction between a gap that target-side observations close and a gap that they do not would fail if any gap yielded to both remedies. Here is a gap of the second kind that a representation closes. Bootstrap over out-of-fold predictions, resampled by stimulus, puts the happiness gain at +0.132 (95\% CI [+0.059, +0.209]), its excess over the mean of the other seven axes at +0.087 ([+0.018, +0.159]) and the tonal-over-spectro- temporal margin on happiness at +0.106 ([+0.007, +0.208]); the last two lower bounds sit close to zero, and because the bootstrap resamples predictions without refitting the models the intervals are narrower than a full resampling would give. Six of six algorithms show a larger gain on happiness than the mean of the other axes. Tonality is not the mechanism of the domain barrier: adding tonal descriptors lowers music-to-environmental transfer by 0.021 (18 paired comparisons, p = 0.038 uncorrected), in the predicted direction but about 4\% of a barrier that runs from 0.592 to 0.054. Interval-consonance weights are fixed a priori from the standard ordering of sensory consonance and are not fitted, so they add no degrees of freedom. All audio is loudness-normalised to $-$23 LUFS; cross-validation is GroupKFold-5 and each cell reports the best of six algorithms. Source data are provided as a Source Data file.}

\newcommand{\legendfigfive}{%
\textbf{An acoustic edit moves both domain models the same way, but the effect is reliable per excerpt for only one of three interventions and only at the level of a descriptor dimension.} \textbf{a} Dose--response curves for three interventions applied to 80 source excerpts (40 per domain) at four dose levels, with dose 0 the unmodified control. Curves show the mean predicted probability of the low-arousal class under a model trained on the music corpus and a model trained on the environmental corpus; dotted lines mark the control level. The modelled quantity throughout is P(low arousal), so a rising curve means the edit made the excerpt read as \textbf{calmer}. Slowing spectral change raises the low-arousal probability in both domains; injecting transients lowers it in both. Envelope compression is retained as a failed manipulation. The curve is an average over 80 excerpts and cannot by itself distinguish an effect from a mean shift; \textbf{d} does that. \textbf{b} Manipulation check, run before any dose--response was read. Within-clip Spearman $\rho$ between dose and the targeted descriptor; point size gives the fraction of clips moving in the same direction as the dose, and the shaded band marks |$\rho$| < 0.5. A target descriptor counts as moved only if |$\rho$| $\ge$ 0.5 and at least 75\% of clips move in the same direction, a criterion fixed in advance. Transient injection meets it on both of its target descriptors and spectral smoothing on two of three; envelope compression meets it on one of three and is carried through the figure as a failed manipulation rather than as a null effect. \textbf{c} Specificity diagnostic. The five descriptors moved most by each intervention at maximum dose, in units of the training-set standard deviation. \textbf{d} Significance against reliability, per intervention and per domain model. Abscissa, the fraction of the 80 excerpts whose own slope across the four doses has the intended sign; ordinate, $-$log10 of a sign-flip permutation p on those 80 slopes. The dashed line at 50\% is chance, and it is the line that matters: a point near it describes an average over excerpts that disagree, not an effect the intervention has on audio in general, however small its p value. Filled markers meet the pre-declared reliability criterion --- sign consistency $\ge$ 65\% and p\_perm < 0.01 --- and open markers do not. Spectral smoothing meets it in both domains (median slope +0.015 music, +0.074 environmental; sign consistency 70\% and 71\%; p\_perm = 1 $\times$ $10^{-4}$ in both). Transient injection reaches significance (p\_perm = 1 $\times$ $10^{-3}$ music, 8 $\times$ $10^{-4}$ environmental) at 50\% and 59\% sign consistency, at or barely above chance, and is therefore reported as a mean shift with no reliable per-excerpt effect and is not built upon. Envelope compression is uninterpretable here, having failed \textbf{b} (p\_perm = 0.50 and 0.98). The directions in \textbf{a} and \textbf{d} cross the domain boundary that Fig. 1 shows the predictions cannot: a model fitted to music does not know how loud is loud in an environmental recording, but it still ranks smoother as calmer, so what fails at that boundary is calibration rather than sign. Panel \textbf{c} bounds what may be claimed: no intervention moves its target descriptor most. Transient injection moves \texttt{spec\_contrast6\_dmean} by 9.0 SD while moving \texttt{onset\_rate} by 2.0 SD, ranking the target twelfth. Because the co-moving descriptors are themselves temporal-variability measures, these experiments support a claim about the variability *dimension* and not about any individual descriptor, and they are statements about model behaviour: no listener response to these edits was measured. An earlier version of this analysis averaged the 80 excerpts within dose and correlated the four resulting points; at n = 4 the smallest attainable two-sided p for a Spearman correlation is 0.083, so that statistic could not have reached significance whatever the data showed, and the excerpt-level test in \textbf{d} replaces it. All interventions are re-normalised to $-$23 LUFS afterwards so that loudness cannot drive the effect. Source data are provided as a Source Data file.}

\newcommand{\legendfigsix}{%
\textbf{Reference distributions, positive controls, and the comparison they license.} \textbf{a} Null distribution of the number of descriptors shared by the top-20 SHAP sets of three models under 20 label permutations. The observed value (4) falls at the null mean (3.75), p = 0.381. Shuffled labels produce the same overlap as intact ones, so the statistic cannot separate signal from noise and no claim is built on it. \textbf{b} Cross-validated AUC under the same permutations, confirming that permutation destroyed the signal (null mean 0.496) while the intact labels reach 0.936. The failure in \textbf{a} is therefore a property of the statistic, not of the data. \textbf{c} Positive controls for two physiological pipelines addressing the same question. In BIRAFFE2 the habituation control is significant in the direction opposite to the established effect ($\rho$ = +0.113 across trials, p = 0.002; skin conductance responses are expected to decline), and post-stimulus windows do not exceed random windows (p = 0.267); applied per participant, none of 94 passed both controls and habituation ran in the expected direction for only 35\%, so no usable subset exists and that pipeline is withdrawn rather than reported. In ds002721 all three controls pass: relative occipital alpha exceeds frontal alpha in 31 of 31 participants (p = 1.9 $\times$ $10^{-9}$); blink-locked averages reach 164 \uV{} frontally, 3.98-fold the occipital amplitude and far above random windows (p = 9.3 $\times$ $10^{-10}$); and a sound-onset response is present (\textbf{d}). Dashed line, p = 0.05. \textbf{d} Grand-average response at fronto-central electrodes (F3, Fz, F4, C3, Cz, C4), time-locked to sound onset, n = 31. Dark band, s.e.m.; grey band, the 95th percentile of a subject-wise sign-flip null with max-statistic correction across the epoch; green shading, the interval exceeding it (372--588 ms). Peak |t| = 8.78 at 452 ms, p = 0.0001 (10,000 permutations). The pre-stimulus interval is not significant (p = 0.54), excluding a slow drift. No classical N1 is present; the stimuli have sharp onsets (median 25 ms to half-amplitude, Extended Data), so the absent fast component is attributed to playback-latency jitter, which limits millisecond-scale analysis but not trial-wise spectral measures. \textbf{e} Stimulus-level intraclass correlation for eight self-report scales and 156 EEG measures, computed on the same 31 listeners, the same 1,240 trials and with the same estimator after within-listener z-scoring. Points are individual measures, vertical bars the median. The EEG set spans four families: regional band power (26), magnitude-squared coherence between region pairs (50), imaginary coherency (50, robust to volume conduction) and electrode-pair asymmetry (30); coherence and asymmetry are included because the original analysis of these recordings located its effect there, and restricting the set to band power would test the conclusion on a representation that excludes the effect in question. Coherence is the strongest EEG family (maximum intraclass correlation 0.092, 95\% CI [0.037, 0.142], split-half 0.51) and band power the weakest (0.040, split-half 0.17), against 0.221 [0.155, 0.281] for the best-agreeing self-report scale on the same trials --- a 2.4-fold gap between two intervals that do not overlap. The figure states one requirement in two settings: a statistic is uninterpretable without the distribution it would take under the null (\textbf{a}, \textbf{b}), and a negative result is uninterpretable until the instrument is shown to detect a known effect (\textbf{c}, \textbf{d}). Panel \textbf{e} is the comparison those controls license, and it is reported as a bound rather than as an absence. The EEG value is the maximum of 156 measures and is biased upward for that reason, so it carries a null built the same way: permuting stimulus labels, recomputing all 156 intraclass correlations and taking the largest, 400 times, places the observed value at p = 0.02, while Benjamini--Hochberg across the same 156 does not pass it (smallest q = 0.16). The two corrections disagree because the value sits between them --- at rank one the Benjamini--Hochberg threshold is the Bonferroni threshold --- and 0.092 is close to the 0.10 this design resolves with 80\% power, so the interval is consistent with values from near zero up to about a tenth of the response variance. What the data support is the comparison and not a zero: whatever stimulus-specific structure the EEG carries is less than half of what the listeners' own ratings carry on the same trials. This is a statement about the component shared across listeners, not about whether the brain responds to sound --- the onset response in \textbf{d} is large and highly reliable; what is weak is the differentiating component. An earlier version of panels \textbf{c}--\textbf{e} was built on an electrodermal analysis that was retracted when its pipeline failed the controls now shown in \textbf{c}, and an earlier stimulus-level value of 0.892 was a leak --- the stimulus identifier itself had entered the measure set --- and is withdrawn. Both are documented in the Supplementary Information. Source data are provided as a Source Data file.}

\newcommand{\legendfigseven}{%
\textbf{The individual-level question is a design problem rather than a null result, and its price depends on the setting.} \textbf{a} Detection measured in real physiology. In an independent corpus of 75 listeners with about 240 auditory events each, the evoked response to a tone was tested per person by a sign-flip permutation test on the difference between real cues and randomly placed anchors. Each point gives the proportion of listeners in whom the response was detected when only n of that listener's observations were used (40 resamples per point). Three channels are shown, each with the analytic power curve for its measured effect size overlaid as a pale band; empirical and analytic differ by 0.015 on average. The dotted line is the false-positive rate obtained by sign-randomising the same differences, which preserves the noise and removes the effect; it stays at nominal (0.03--0.07) across the whole range, so the rise in the curves is detection and not false positives. The shaded band marks what a laboratory session supplies. Within it, an effect plainly visible in the group average ($-$5.6 \uV{} at the N1, group p = 0.0002 in all three channels) is recovered in 53\% of individuals from EEG, 39\% from pupil diameter and 23\% from heart rate. \textbf{b} Probability of detecting a within-person association as a function of the observations available per person, for four true association strengths. Synthetic data were generated from the empirical feature covariance of the EEG corpus (Ledoit--Wolf shrinkage) with an association of known strength injected, then passed through the identical pipeline, criterion and per-person permutation test used for the real data (500 simulations per cell, 100 permutations each, n from 30 to 5,000). Within the laboratory-session band, power is 0.10 for r = 0.40, 0.072 for r = 0.30, 0.06 for r = 0.20 and 0.04 for r = 0.10. The r = 0.10 curve is not monotone below n = 250; the Monte-Carlo s.e. is 0.009 to 0.015 per cell and the curve is included for the floor it establishes, not for its shape. \textbf{c} Observations per person required for 80\% detection, measured and simulated on one axis, against three supply levels. Every requirement is resolved by this grid; the two weakest effects, which a coarser run could only bound at > 1,000, now stand at 1,364 (r = 0.20) and 4,888 (r = 0.10). Dotted lines give the supply from this corpus (30), comparable public corpora (60), and a single overnight recording sampled once per minute over eight hours (480). Heart rate, the signal a consumer device records, requires roughly four times as many observations as scalp EEG (233 against 57). \textbf{d} When per-person calibration begins to pay, in three settings. For each descriptor--outcome cell the between-person standard deviation of the acoustic slope was estimated by restricted maximum likelihood (filled circle, median over that setting's cells) and compared with a null built by parametric bootstrap from that cell's own design (cross, median over the same cells). The null is plotted rather than subtracted, because a between-person spread estimated from few observations per person is biased upward: at 18 observations per person the estimator returns 0.11 when the true value is zero, and a reader who cannot see that cannot judge the physiological row. Break-even, n*, is the number of observations per person at which the corrected spread exceeds the error in estimating it. Urban soundscape ratings (42 observations per person, 8 cells): estimate 0.134 against a null of 0.013, n* = 67, and 2 of 8 cells have already crossed break-even, both on the eventfulness axis. Music ratings (44 per person, 10 cells): 0.074 against 0.000, n* = 166, none across. Music electrodermal responses (18 per person, 15 cells): 0.106 against a null of 0.114 --- the estimate does not exceed its own noise --- with n* = 318 taken over the 5 cells in which it is finite, 11 of the 15 indistinguishable from their own null, and none across. Both kinds of estimate in \textbf{a}--\textbf{c} are lower bounds, for different reasons. The simulation assumes multivariate normality, a linear association and within-person stationarity, each of which favours detection. The measurement uses a deliberately plain pipeline --- a fixed-window mean over one electrode cluster, no independent component analysis and no spatial filtering --- so a dedicated pipeline would do better; it is, however, representative of what a non-specialist pipeline achieves. The extrapolation behind \textbf{c} was validated out of sample: predicting detection at n from an effect size estimated on a disjoint half of the same person's trials matches the observed rate to within 0.03 on average. Panel \textbf{d} is a different quantity from \textbf{a}--\textbf{c} --- break-even for per-person calibration, not detection of a within-person association --- and it is the one that varies, five-fold across settings. Two earlier claims from that analysis are withdrawn and do not appear here: a between-person spread of 0.106 for the physiological setting, which sat below its own null; and the reading of a confidence interval excluding zero as evidence of individual differences, since that interval is for the composite of true spread and null and, at 18 observations per person with a true spread of zero, excludes zero on 84\% of occasions. Read together, \textbf{a}--\textbf{c} place the shortfall at roughly a factor of twenty and locate it in duration and repetition rather than in sample size: the quantity in short supply is observations within a person, which recruiting more participants does not provide. Source data are provided as a Source Data file.}

\newcommand{\legendfigeight}{%
\textbf{The four boundaries divide into two kinds of failure with mutually exclusive remedies.} \textbf{a} What survives each crossing, as a fraction of the ceiling attainable on that target, ordered near-to-far --- the order in which an application meets them. Colour carries the diagnosis rather than the boundary: budget-limited, where target-side observations close the gap, and information-limited, where they do not. Cross-corpus contributes two bars because it is the only boundary on which both kinds appear inside one experiment, with the same models, descriptors and act of swapping one corpus for another: within a domain 0.614 of 0.769 survives (80\%), across domains 0.047 of 0.769 (6\%). Cross-synthesis is drawn at 1 and 0 by construction and is not a measured fraction; the outcome there is categorical, in that the sign of an acoustic edit crosses the domain boundary while the calibration does not. Cross-channel is the best of 156 EEG measures against the best-agreeing self-report scale on the same 1,240 trials (0.092 of 0.221, 42\%). Cross-individual is the observations a corpus supplies against those required for per-person calibration to pay, shown at the two ends of its five-fold range (soundscape ratings 42 of 67, 63\%; music electrodermal responses 42 of 318, 13\%). \textbf{b} What it costs to buy the gap back, on one axis of target-side observations. Filled bars are prices that exist: 100 labelled target examples for two-thirds of the gap a corpus swap opens (Fig. 2, where that price is measured on pairs that cross the music/environmental boundary), and 67, 166 and 318 observations per person for break-even on per-person calibration in urban soundscape ratings, music ratings and music electrodermal responses (Fig. 7d). Three rows carry no bar but an open marker at the right edge with the reason. The marker reads "no price found" and means exactly that: under the remedies tested, no purchase recovered the gap. It is not a claim that none exists at any budget, and the distinction matters because only the priced rows rest on a target-side budget curve. No representation crossed the domain boundary --- four pretrained representations spanning supervised, self-supervised and contrastive pretraining, none of them recovering it (Fig. 1f). The calibration deficit at the synthesis boundary is not a shortage of data, so no quantity of it is the remedy. At the response-channel boundary no information source exceeds 31\% of the attainable ceiling on the physiological target, each route scored against the ceiling of the electrodermal measure its correlation was computed on; three routes stopping within a third of it makes it unlikely that the limit lies in the predictor, but the two acoustic routes sit only just above a random projection of the same descriptor space (21--25\% against 21\%), and no target-side learning curve was measured on that boundary, so saturation was not demonstrated. A price that was not found is drawn as an open marker rather than as a very long bar, because a long bar would assert that the purchase exists and is merely expensive. \textbf{c} The pre-declared falsification test. The division in \textbf{a} is worth nothing if it cannot fail, and the condition was fixed in advance: a gap closable both by additional observations and by a change of representation would break it. Gain in cross-validated Spearman $\rho$ from adding 39 tonal descriptors, per affect axis of the Soundtracks corpus (n = 360). The criterion --- that the gain concentrate on happiness, at more than twice the mean gain of the other seven axes and above 0.05 --- is met at +0.132 against +0.045 (Fig. 4). This figure re-plots quantities established elsewhere and introduces no new analysis; every value traces to the figure cited beside it and to the same Source Data. The three quantities are defined per boundary and are not interchangeable between them: the denominator in \textbf{a} is the within-corpus performance of the same representation for boundary 1, the within-corpus performance of the model being perturbed for boundary 2, the ratings' stimulus-level reliability on the same trials for boundary 3, and the observations required for break-even for boundary 4. Source data are provided as a Source Data file.}

\newcommand{\legendedone}{%
\textbf{Loudness acts as a corpus fingerprint rather than as signal.} \textbf{a} Cross-corpus AUC (DEAM $\rightarrow$ PMEmo) for five classifiers before and after loudness normalisation to $-$23 LUFS. Normalisation raises transfer for every model, from a mean of 0.866 to 0.886. \textbf{b} Mean SHAP rank of the five level-dependent descriptors before and after normalisation; arrows run from raw to normalised. All five lose importance, with \texttt{rms\_p90} falling from rank 39 to rank 93. Removing absolute level improves rather than degrades cross-corpus transfer, which is the signature of a corpus-specific confound: mastering practice differs systematically between corpora. Because every transfer number in this work is computed on normalised audio, this control is the admission criterion for the domain barrier, and the obvious objection --- that normalisation creates the barrier by discarding a genuinely informative dimension --- was tested separately. Three feature pipelines that share no data, no intermediate table, no constant and no estimate --- normalised descriptors, descriptors computed with loudness retained, and the 6,373-dimensional ComParE set extracted from the original recordings by the corpus authors --- put the best cross-domain transfer at $\rho$ = 0.045, 0.049 and 0.072. The spread of 0.027 is below the 0.05 threshold fixed before the test, and loudness alone contributes only +0.059. This is the only result in this work for which agreement between separately derived pipelines is claimed; no other comparison here meets that condition. Source data are provided as a Source Data file.}

\newcommand{\legendedtwo}{%
\textbf{Regression framing recovers the discarded middle of the arousal range.} \textbf{a} Clips available, used under the regression framing, and used under the original binary framing, per corpus. \textbf{b} Totals across corpora. The binary framing selects the two extremes of the arousal axis and discards the middle, retaining 1,341 of 3,782 annotated clips. The regression framing uses 3,163. The discarded clips carry the threshold information that a rule library requires. Source data are provided as a Source Data file.}

\newcommand{\legendedthree}{%
\textbf{Glass-box models cost almost no accuracy.} \textbf{a} Spearman $\rho$ under GroupKFold-5 for black-box and glass-box models in both domains. \textbf{b} Change in $\rho$ from imposing monotone constraints on the seven cross-domain-invariant descriptors, relative to the same feature subset without constraints. Monotone constraints cost +0.0003 in the ambient domain and $-$0.0016 in the music domain --- replicated at zero cost in both. The monotonicity assumed by threshold rules is therefore supported by the data rather than imposed on it. The fully additive EBM trails the best black box by 0.019. Source data are provided as a Source Data file.}

\newcommand{\legendedfour}{%
\textbf{Explainable Boosting Machine shape functions give thresholds without post-hoc attribution.} Additive contribution to predicted arousal as a function of descriptor value for the six most important terms in the ambient-domain EBM (n = 1,213, interactions disabled). Red lines mark the zero crossing, printed above each panel; these are the thresholds a rule library reads directly, with no attribution step. Five of the six terms are temporal-variability measures. Source data are provided as a Source Data file.}

\newcommand{\legendedfive}{%
\textbf{Descriptor importance is stable across the Rashomon set.} Share of near-optimal models ranking each descriptor in their top 15. The Rashomon set comprises the 10 of 60 sampled models within 0.02 Spearman $\rho$ of the best ($\rho$ $\in$ [0.823, 0.843]); models were sampled over algorithm family, depth, regularisation and subsampling, and importance was measured by permutation importance on a held-out fold. Six descriptors are recovered by every near-optimal model, four of them temporal-variability measures. Reporting importance from a single fitted model would have left this multiplicity untested. Source data are provided as a Source Data file.}

\newcommand{\legendedsix}{%
\textbf{No EEG measure tracks self-reported emotion across listeners at the effect size this design can resolve.} Each point is one of 208 combinations of an EEG measure (26) and a rating scale (8). The abscissa gives the mean across the 31 listeners of the within-listener partial Spearman $\rho$ between that measure and that listener's own trial-by-trial rating, with trial index partialled out; both EEG and ratings can drift over a session, and a shared drift would otherwise manufacture a correlation. The ordinate gives the evidence against a zero group mean (Wilcoxon signed-rank). Point size is the fraction of listeners whose sign matches the group mean; colour is the rating scale. The dashed line marks the Benjamini--Hochberg threshold at FDR < 0.05. No combination crosses it (largest |mean $\rho$| = 0.115; median 0.028; the group mean's sign is shared by a median of 55\% of listeners), and the distribution is symmetric about zero as expected under the null. This is a bound and not an absence. Listeners contributed a median of 30 rated trials each, and the design simulation (Fig. 7b) shows that at 30 observations per person a within-person association of r = 0.40 would be detected on 10\% of occasions and one of r = 0.20 on 6\% --- that is, at chance. The panel therefore establishes only that no large within-person association is obvious, and supports no conclusion about associations below approximately r = 0.5. The group-level result in Fig. 6e does not share this limitation, because it pools 31 listeners at the level of the stimulus. This figure is computed on the pipeline validated in Fig. 6c,d, which is what makes the negative interpretable at all; independent-component removal of ocular artefacts, used here as a sensitivity manipulation check, increased detection of the known effect while further reducing the effect of interest. An earlier version used an electrodermal analysis that was retracted when its pipeline failed the same controls. Source data are provided as a Source Data file.}

\newcommand{\legendedseven}{%
\textbf{A pretrained representation gains most on the one corpus that overlaps its pretraining data.} Best cross-validated Spearman $\rho$ for 122 hand-crafted spectro-temporal descriptors and for 512-dimensional pretrained audio embeddings, on three corpora (best of six algorithms, GroupKFold-5, loudness-normalised audio). Embeddings win on all three: +0.040 on DEAM (0.736 $\rightarrow$ 0.776), +0.034 on Soundtracks (0.823 $\rightarrow$ 0.856) and +0.070 on Emo-Soundscapes (0.773 $\rightarrow$ 0.843). The advantage does not require a strong downstream model --- on embeddings, cross-validated linear ridge regression is competitive with the tuned tree ensembles --- so it is carried by the representation and not by the estimator. The two music corpora have no overlap with the embedding model's pretraining data and give the two smaller margins; the corpus that derives from the public sound repository appearing in that pretraining data gives the largest. We report the music-corpus figure, +0.037, as the estimate that generalises and leave the difference unattributed. It cannot be assigned to pretraining overlap here: the one corpus sharing data with pretraining is also the only environmental corpus in the comparison, and both clean corpora are music, so overlap and sound domain are collinear. The larger margin is equally consistent with the embedding model simply being stronger on environmental sound. Separating the two accounts would require an environmental corpus outside the pretraining data, or a music corpus inside it. This figure concerns within-corpus accuracy only. Whether a pretrained representation crosses the domain boundary is a separate question, answered for four pretraining paradigms in Fig. 1f and Table 1, where none of them does. Source data are provided as a Source Data file.}

\newcommand{\legendedeight}{%
\textbf{Only temporal-variability descriptors keep their direction across sound domains.} \textbf{a} Standardised effect size (Cohen's d) of each FDR-significant descriptor in the music domain versus the environmental domain, with marginal densities by feature class. Shaded quadrants mark sign preservation; filled circles keep their sign across domains, open circles reverse it. The dotted diagonal is the identity line. Descriptors with min|d| > 0.35 are labelled. \textbf{b} The eight sign-preserving descriptors, ranked by the weaker of their two effect sizes (printed at the right of each row). Circles give |d| in each domain and the connecting bar gives the gap between them. All eight are negative in both domains: lower descriptor value corresponds to lower arousal, whether the sound is music or an environmental recording. The descriptors that reverse are without exception measures of absolute spectral level --- the property on which the two domains differ most. Feature classes are assigned by statistic type --- descriptors ending in \texttt{\_std} or \texttt{\_dmean}, plus \texttt{onset\_rate}, are temporal-variability measures; all others are spectral-level measures. Significance was established per domain by per-feature permutation testing (200 permutations) with Benjamini--Hochberg control at q = 0.10. This analysis is the second of the four selection routes in Fig. 3, and is the only one that tests the boundary directly. Source data are provided as a Source Data file. This analysis appeared as a main figure in an earlier version; it was moved to Extended Data when the price of the first boundary (Fig. 2) took the main-figure slot.}
